% =========================================================================
% GUÍA DE DEFENSA
% =========================================================================
\chapter*{Anexo: Guía de Defensa}
\addcontentsline{toc}{chapter}{Anexo: Guía de Defensa}

\section*{Preguntas Generales}
\textbf{1. ¿Cuál es la contribución central de esta tesis?} Un pipeline reproducible que integra econometría panel y aprendizaje automático para predecir el deterioro financiero de SOFIPOs mexicanas, con interpretabilidad SHAP y validación temporal rigurosa.

\textbf{2. ¿Por qué SOFIPOs y no bancos comerciales?} Porque atienden poblaciones más vulnerables a shocks, cuentan con menor sofisticación analítica interna, y presentan mayor heterogeneidad estructural, haciendo más valiosa una herramienta panel de alerta temprana.

\section*{Preguntas sobre Datos y Metodología}
\textbf{3. ¿Cuántas instituciones y observaciones se usaron realmente?} 126 instituciones y 1690 observaciones en el panel bruto; 46 instituciones y 1189 observaciones en el panel final de modelado, tras exigir historia mínima de 6 trimestres.

\textbf{4. ¿Por qué la partición train/test es cronológica y no aleatoria?} Porque una partición aleatoria en series temporales permitiría que el modelo "vea" información futura durante el entrenamiento (\textit{lookahead bias}), invalidando la evaluación.

\textbf{5. ¿Qué es \texttt{TimeSeriesSplit} y por qué se usó?} Un esquema de validación cruzada que en cada partición entrena con el pasado y valida con un bloque futuro inmediato, preservando el orden temporal, a diferencia del \textit{k-fold} estándar.

\textbf{6. ¿Por qué el NICAP promedio reportado en el boletín crudo es de 88{,}627\%?} Por valores atípicos extremos en un número reducido de observaciones (hasta $4.97\times10^7\%$), documentados explícitamente como anomalía de calidad de datos, no corregidos arbitrariamente sino señalados y tratados mediante estadística robusta en el modelado.

\section*{Preguntas sobre Modelos}
\textbf{7. ¿Por qué XGBoost y no solo regresión lineal?} Porque la relación entre las variables predictoras y el IMOR futuro presenta interacciones no lineales (confirmadas por la importancia dominante de variables no lineales como medias móviles), que un modelo lineal no captura adecuadamente.

\textbf{8. ¿Qué significa un $R^2$ de 0.6126?} Que el modelo explica el 61.26\% de la varianza del IMOR futuro en el conjunto de prueba fuera de muestra; el 38.74\% restante corresponde a factores no capturados por las variables disponibles (shocks idiosincráticos, variables omitidas).

\textbf{9. ¿Qué significa un MAE de 3.29\%?} Que, en promedio, la predicción del IMOR futuro se desvía 3.29 puntos porcentuales del valor real observado.

\textbf{10. ¿Por qué ARIMA tiene menor MAE que XGBoost?} Porque se evalúa sobre un subconjunto muy reducido y favorable (3 instituciones con la historia más extensa y estable), mientras que XGBoost se evalúa sobre el panel completo de prueba, mucho más heterogéneo; no son comparables en igualdad de condiciones muestrales.

\textbf{11. ¿Por qué SMOTE-ENN si el desbalance era solo del 49\% al 51\%?} Porque, aunque el desbalance era moderado y no severo, SMOTE-ENN también cumple una función de depuración de observaciones ambiguas en la frontera de decisión (fase ENN), mejorando la separabilidad de clases más allá del simple balanceo de proporciones.

\textbf{12. ¿Por qué ROC-AUC y MCC como métricas principales de clasificación?} Porque la exactitud simple puede ser engañosa ante cualquier grado de desbalance; ROC-AUC mide la capacidad de discriminación independientemente del umbral, y MCC pondera de forma simétrica los cuatro cuadrantes de la matriz de confusión.

\textbf{13. ¿SHAP demuestra causalidad?} No. SHAP cuantifica la contribución marginal de cada variable a la predicción \textbf{dentro del modelo entrenado}; refleja asociación estadística aprendida, no un mecanismo causal verificado en el sistema económico real.

\textbf{14. ¿Qué pasa si una SOFIPO nueva no tiene historial?} El pipeline exige un mínimo de 6 observaciones trimestrales para calcular consistentemente los rezagos de orden 3; instituciones más nuevas quedan fuera del panel de modelado, una limitación explícita reconocida en el Capítulo 1.

\textbf{15. ¿Qué es \textit{data leakage} y se descartó en este trabajo?} Es la contaminación del entrenamiento con información no disponible en el momento de la predicción. Se descartó revisando que: la partición es cronológica, SMOTE-ENN se aplica solo post-partición sobre entrenamiment, y los targets se construyen por desplazamiento hacia adelante dentro de cada institución (Capítulo 7).

\textbf{16. ¿Qué haría diferente con más datos?} Sustituir las series macroeconómicas sintéticas por series oficiales de Banxico, extender el horizonte de predicción, y evaluar arquitecturas LSTM con suficiente profundidad temporal por institución.

\textbf{17. ¿Qué ocurriría durante una crisis sistémica no vista en el entrenamiento?} El desempeño no puede garantizarse: los modelos de aprendizaje automático extrapolan pobremente fuera de la distribución de entrenamiento; se recomienda \textit{stress testing} explícito como línea de trabajo futuro, no realizado en esta tesis.

